
\section{Introduction}


Topography on the Earth's core mantle boundary (CMB) likely plays an important role in the dynamics of the core, with possible observable effects in both the geomagnetic field \citep{aA77} and changes in the length of day ($\Delta$LOD) \citep{dJ89,HdV05,pR12,fM25,sR23}. 
Stationary geomagnetic flux patches and initiation sites of magnetic polarity reversals are hypothesized effects of CMB topography \citep{jT15}.
Flow in the core will lead to pressure variations along the topographic surface, resulting in torques on the mantle and exchanges of angular momentum that contribute to $\Delta$LOD \citep{rH69,wK97,pR12}. 
These topographic torques are not the only source of $\Delta$LOD. Gravitational coupling between the inner core and the mantle \citep{bB96,jM03} and electromagnetic coupling between core-generated magnetic field and electrically conducting regions at the base of the mantle \citep{bB07,sB70, dJ15,wK19} both likely play a role in generating $\Delta$LOD. Nevertheless, recent work has demonstrated that topographic torques exhibit scalings that indicate they are of sufficient magnitude to play a significant role in $\Delta$LOD \citep{tO25}. However, this work was limited to a spatially localized topography with a fixed width, whereas global scale CMB topography is likely present. Prior studies investigating the ramifications of CMB topography have been limited to the linear regime \citep{pB96}, or simplified two dimensional models that may not fully capture all dynamical effects \citep{jH98,mW03,mC12}. 
The primary goal of the present work is to determine how globally distributed topography influences both the topographic torques and core dynamics using a suite of fully three dimensional, turbulent numerical models. 

While the shape of the CMB is poorly constrained, seismic investigations and geodynamic modeling of the deep mantle point to possible structure on the CMB. Seismological studies show considerable variation in inferred CMB structure, with typical topographic amplitudes up to a few kilometers for global scale wavelengths \citep{kK03,sT10}. An obvious possible cause for CMB topography is mantle dynamics. It has been proposed that regions of slow seismic velocity deep in the mantle, the so-called Large Low Shear Velocity Provinces (LLSVPs) \citep{sC16,pK21}, contribute to the CMB deformation. The largest LLSVPs are located at the southern tip of Africa and in the southern Pacific. These findings suggest that CMB topography may consist of a few large structures, however dynamic mantle modeling has shown that the presence of LLSVPs can generate short wavelength topographic features depending on the thermal and compositional characteristics \citep{tL07,tL10}. Furthermore, the edges of the LLSVPs may be quite sharp, which may imply lateral variations in topography over tens to hundreds of kilometers \citep{sC16}. 

%Geodynamic modeling 
%Both seismic imaging and mantle convection studies suggest that CMB topography has a large lateral extent \citep{kK03,sT10,pK21,tL10}
%
% and may be characterized by the largest amplitudes in comparison to the small scale topography, as indicated by , for instance, by Kaula's rule \citep{wK66,mP23}


% and to the author's knowledge only a single study has considered convection in a sphere with globally distributed topography \citep{wK01}, although numerical constraints precluded turbulent flows. 
%Significant uncertainty exists regarding the shape of CMB topography, so it is useful to understand the effects of varying topographic shape. It has been proposed that regions of slow seismic velocity deep in the mantle, the so-called Large Low Shear Velocity Provinces (LLSVPs), contribute to the CMB deformation \citep{sC16,jT15}. The largest LLSVPs are located at the southern tip of Africa and in the southern Pacific. This suggests that the primary feature of CMB topography is a few large structures, however dynamic mantle modeling has shown that the presence of LLSVPs can generate short wavelength topographic features \citep{tL07}. Furthermore, the edges of the LLSVPs may be quite sharp, which may imply lateral variations in topography over 10's to 100's of kilometers \citep{sC16}. 
%In this study, we investigate a ``tomography model" based on the locations of the LLSVPs. 

The canonical model of core dynamics consists of a rotating, self-gravitating fluid bounded by spherically symmetric surfaces that is heated from within. 
For sufficiently strong heating, the basic instability in this system consists of convective Rossby waves that are characterized by an approximately axially invariant structure with short longitudinal wavelength \citep{pR68,fB70}. 
These waves give rise to broadband turbulence as the system is forced more strongly \citep[e.g.][]{tG16,jN24}, and the resulting motions are known to readily generate magnetic fields that bear similarities with the geomagnetic field for certain parameter values \citep{cF23,juA23,nG22,nG24}. While this system continues to serve as the `zeroth' order model for the geodynamo, it is known that unique dynamical effects arise when the CMB deviates from axial symmetry \citep{cD08,jM23}. 
A well known result for rapidly rotating flows is that the streamlines follow a particular set of geometric curves known as ``geostrophic contours'' \citep{hG68}. 
In a perfect sphere these curves form circles about the origin, and a mechanism to break geostrophy is necessary for the convective instability. In the hydrodynamic case viscosity is the mechanism that breaks geostrophy and it yields lengths scales of size $O\lb Ek^{1/3}\rb,$ where $Ek$ is the Ekman number -- a ratio of viscous to Coriolis forces \citep{sC61}. 
As convection is made more turbulent it is generally found that the convective scale increases from this viscously selected small scale, although uncertainty about the extent of this effect exists \citep{cG19, tO23, jN24, cG25}. 
%t has been suggested that the Rossby number $Ro$- a ratio of inertial to Coriolis forces- determines the convective scales of the flow. 
%The scale $Ro^{1/2}$ is proposed , although it may be an over-estimate \citep{tO23}. 

Bell and Soward (1996) \citep{pB96} and Bassom and Soward (1996) \citep{aB96} showed that when boundary topography is considered, the most unstable modes are characterized by the length scales present in the topography. The convective modes are no longer constrained by viscosity and do not need to be small.
The results of these two studies hinge on the fact that topography breaks the axisymmetry of the geostrophic contours and allows buoyancy to directly drive the flows.
Deviations from rotational symmetry are likely present on the CMB, implying that the geostrophic contours may be characterized by complex shapes. Indeed, observations show that large scale flows in the core are strongly asymmetric \citep{rH11,pL17,cF12,cF23}, indicating that the CMB could be important for core dynamics. Resonance between the flow and topography, for example, can generate time-independent flows which likely have consequences for the structure and dynamics of the geomagnetic field \citep{mC12}.


%In a geophysical context we may expect the topographic length scales to differ greatly from the small scales typically considered in convection studies.  
%In the core, estimates for the lateral extent of CMB topography suggest scales that range from global scales to $O\lb 10 \text{km}\rb$ \citep{aT05}. Even the $Ro^{1/2}$ scaling suggests convective scales up to only $O\lb 10\text{km}\rb $, so it is likely that the topographic and typically studied convective scales are distinct.

%The topographic studies of \cite{pB96} (hereafter referred to as BS96) and \cite{aB96} used an annulus geometry originally pioneered by Busse \citep{fB70}, and a number of results \citep{jH98,mW03,mC12} have demonstrated that sufficiently large topography reduces the critical Rayleigh number in this geometry. So far, however, the application to a sphere has not been explored.
%The $Ek^{1/3}$ and $Ro^{1/2}$ scalings suggest convective scales of $O\lb 10 \text{m}\rb $ and $O\lb 10 \text{km}\rb $  respectively, less than all but the smallest topography scales.


%In a spherical shell the geostrophic contours are concentric rings centered about the origin, however boundary topography serves to deform the contours.
%In general, the contours are no longer circular, nor are they required to form closed loops. 

%Topography can also control the turbulent mixing, leading to enhanced transport in heat and momentum. Resonance between the flow and topography can generate time-independent flows which likely have consequences for the geomagnetic field \citep{mC12}. However, the full extent of the dynamical influences of CMB topography on core flows have yet to revealed given the limited number of investigations on this topic. The present study focuses on improving understanding of how the topographic shape influences these dynamics. 


The present work studies the influence of CMB topography on core convection and the resulting torques through direct numerical simulations that allow for a non-spherical boundary. 
%We perform turbulent simulations for a variety of topographic shapes at fixed rotation rate and thermal forcing. 
The paper is organized as follows. 
In Section \ref{s:methods} we introduce the governing equations for rotating convection, discuss the implementation of CMB topography, and define relevant quantities for the investigation.
Results are given in Section \ref{s:results}, where we discuss heat and momentum transport, flow morphology, and topographic torques. To compare with the theoretical predictions of Bell and Soward (1996) \citep{pB96}, we first present results from a set of simulations in which the thermal forcing is weak. 
%We discuss the instability presented inand the impact of the deformed geostrophic contours. 
The results from the turbulent simulations are then presented, which are all performed at the same rotation rate and thermal forcing.
Five different topographic shapes are considered, including four that are represented by single spherical harmonics, and a fifth topographic shape that is informed by studies of LLSVPs. 
For each topographic shape we vary the topographic amplitude.
We discuss the role of shape and amplitude on the global transport of momentum and heat, as well as the topographic torques exerted on the CMB. 
We demonstrate that the deformation of the geostrophic contours generally serves to increase the heat and momentum transport, but has a more complicated effect on the torques for certain parameters. 
We provide further evidence and justification for the dynamical scaling law which suggests a linear dependence on topographic amplitude and a quadratic dependence on flow speed \citep{tO25}.
We propose a procedure by which we can determine the magnitude of the torques given a particular topographic shape.
In Section \ref{s:disc} we discuss our results and provide geophysical context. 
